\par\Needspace{6\baselineskip}\medskip
\begin{center}\hypertarget{native-171-article124}{}\pdfbookmark[2]{124.}{article-124}124.\end{center}
We now enumerate the most suitable methods based on the preceding principles. Their main features were set out in Book One and need only be adapted here to our present purpose.

The simplest method seems to be to take as $x,y$ the distances of the celestial body from the Earth at two observations, or rather either the logarithms of those distances or the logarithms of the distances projected onto the ecliptic or the equator. From these we obtain, by Article 64, V, the heliocentric positions and the corresponding distances from the Sun. Then, by Article 110, we find the orientation of the orbital plane and the heliocentric longitudes in the orbit. From these, together with the corresponding radius vectors and times, we obtain all the remaining elements by solving the problem treated at length in Articles 85–105. These elements clearly represent the two observations exactly, whatever values were assigned to $x,y$. If we now use these elements to compute the geocentric position at the time of the third observation, its agreement or disagreement with the observed position will determine whether the assumed values were the true ones or differed from them. This gives two comparisons, and we may take one difference, in longitude or right ascension, as $X$, and the other, in latitude or declination, as $Y$. Unless, therefore, the values of these differences $X,Y$ come out $=0$ of their own accord, the true values of $x,y$ may be found by the method described in Articles 120 and following. In itself the choice of which two of the three observations to start from is arbitrary. It is generally preferable, however, to use the first and last, apart from the special case we shall discuss presently.

\SourceLeaf{172}This method is preferable to most of those to be explained later because it admits of the most general application. An exception must be made when the two extreme observations cover a heliocentric motion of 180, 360, or 540 degrees, and so on. In that case the orientation of the orbital plane cannot be determined from two heliocentric positions (Article 110). It is likewise inadvisable to apply the method whenever the heliocentric motion between the two extreme observations differs very little from $180^0$ or $360^0$, and so on. In such a case the orientation of the orbit cannot be determined accurately. More precisely, the slightest changes in the assumed values of the unknowns would produce such large changes in the orientation of the orbit, and hence also in the values of $X,Y$, that the latter could no longer be regarded as proportional to the former. A remedy is available, however. In such a case we start from the first and middle observations or from the middle and last, taking as $X,Y$ the differences between calculation and observation at the third or first position, respectively. If, however, the second position were about 180 degrees from the first and the third about 180 degrees from the second, the difficulty could not be removed in this way. It is better to leave such observations out of the calculation of the elements, since by the nature of the case they afford no accurate determination at all of the orientation of the orbit.

The method also makes it easy to estimate how much the elements change if the two extreme positions are held fixed and the middle one is varied slightly. In this way we can form some judgment of the degree of precision to be attributed to the elements found.
